
\usepackage[utf8]{inputenc}
\usepackage[T1]{fontenc}
\usepackage{charter}
% \usepackage[bitstream-charter]{mathdesign}
\usepackage[francais,english]{babel}
\usepackage[x11names, rgb, usenames, dvipsnames]{xcolor}
\usepackage{geometry}
\usepackage[pdftex]{graphicx}
\usepackage{tikz}
\usepackage{pgfplots}%Needed at least by {axis} environment
\pgfplotsset{compat = 1.10}%1.3 needed by 'ylabel shift = 1 em'
\usetikzlibrary{patterns}

\usepackage{subfig}

\usepackage{graphicx}
\usepackage{listings}
\usepackage{picture}
\usepackage{epic}
\usepackage{verbatim}
\usepackage{textcomp}
\usepackage{tabularx}
\usepackage{multirow}
\usepackage{array}
% \usepackage[breaklinks=true,hidelinks]{hyperref}
\usepackage{memhfixc}
\usepackage[breaklinks=true]{hyperref}

\usepackage{amsmath}
\usepackage{mathtools} %dcases
\usepackage{dsfont} %Needed for \mathds{N}

\usepackage[numbers]{natbib}

\usepackage{framed}
\usepackage{listings}
\definecolor{colKeys}{rgb}{0,0,1}  % couleurs des mots-clés propres au language
\definecolor{colIdentifier}{rgb}{0,0,0}  % couleurs des mots à identifier
\definecolor{colComments}{rgb}{0,0.6,0}  % couleurs des commentaires
\definecolor{colString}{rgb}{0.6,0.1,0.1}

\lstset{commentstyle=\color{colComments}\textit,
float=hbp,%
basicstyle=\ttfamily\footnotesize, %
identifierstyle=\color{colIdentifier}, %
keywordstyle=\color{colKeys}, %
stringstyle=\color{colString}, %
columns=flexible, %
tabsize=2, %
extendedchars=true, %
showspaces=false, %
showstringspaces=false, %
numbers=left, %
numberstyle=\tiny, %
breaklines=true, %
breakautoindent=true, %
captionpos=b,
frame=lines}

\lstdefinelanguage{diff}{
  morecomment=[f][\color{NavyBlue}]{@@},     % group identifier
  morecomment=[f][\color{BrickRed}]-,         % deleted lines 
  morecomment=[f][\color{ForestGreen}]+,       % added lines
  morecomment=[f][\color{magenta}]{---}, % Diff header lines (must appear after +,-)
  morecomment=[f][\color{magenta}]{+++},
}

\lstdefinelanguage{config}{
  morecomment=[f][\color{MidnightBlue}]\#,%comment
}

%%%%%%%%%%%%%%%%%%%%%%%%%%%%%%%%%%%%%%%%%%%%%%%%%%%%%%%%%%%%%%%%%%%
%http://tex.stackexchange.com/questions/2393/testing-for-current-page
%%%%%%%%%%%%%%%%%%%%%%%%%%%%%%%%%%%%%%%%%%%%%%%%%%%%%%%%%%%%%%%%%%%
\usepackage{ifthen}
\newcounter{cPage}
\newcommand\figureref[1]{%
  figure~\ref{#1}%
 \refstepcounter{cPage}\label{current\thecPage}%
 \ifthenelse{\equal{\pageref{#1}}{\pageref{current\thecPage}}}%
  {}{~(p.~\pageref{#1})}}%
%%%%%%%%%%%%%%%%%%%%%%%%%%%%%%%%%%%%%%%%%%%%%%%%%%%%%%%%%%%%%%%%%%%




\setcounter{secnumdepth}{3}% Number up to subsubsection
\maxsecnumdepth{subsubsection}% 'secnumdepth' is reset by \mainmatter. You should use 'maxsecnumdepth'.
% \seccounter{tocdepth}{3}% Add subsubsection to the table of content.




%%%%%%%%% vvvvvvv
% Anthony Debruyn 
\let\oldCaption\caption
\renewcommand{\caption}[2]{
\oldCaption[#1]{{\small\bfseries #1:} #2}
}
%%%%%%%%%%%%%%%%%%%%%%%%%%%%%%%%%%%%%%%%%%%%%%%%%%%%%%%
% NEW STYLES
%%%%%%%%%%%%%%%%%%%%%%%%%%%%%%%%%%%%%%%%%%%%%%%%%%%%%%%
\captiondelim{ -- }
\captionnamefont{\small\bfseries}
\captiontitlefont{\small}
\precaption{\rule{\linewidth}{0.4pt}\\}
% Anthony Debruyn
%%%%%%%%% ^^^^^^^




% ----------------------
% lists
% ----------------------
\usepackage{enumitem}
% \setlist{nosep} % remove all vertical sep inside and around the list
\setlist{noitemsep} % remove all vertical spaces inside the list
% \setlist[2]{noitemsep} % set the itemsep and parsep for all level two lists to 0
% \setenumerate{noitemsep} % set noitemseponly for enumerate
% \begin{description}[noitemsep] % if the \setlist{noitemsep} as not been defined, use 'noitemsep' only for this list.
% \end{description}


% ----------------------
% Cover page
% ----------------------

% Logos
\newcommand{\ulb}{\includegraphics[scale=1.1]{logo_ULB2.pdf}}
\newcommand{\polytech}{\includegraphics[scale=0.35]{logo_polytech_FR.pdf}}

% Polices
\definecolor{ULBblue}{rgb}{0,0.2196,0.5765}
\newcommand{\fontTitle}{\sffamily \Huge\selectfont \color{ULBblue}}
\newcommand{\fontSubtitle}{\sffamily \LARGE \selectfont \color{ULBblue}}
\newcommand{\fontText}{\sffamily \selectfont}
\newcommand{\fontColor}{\sffamily \selectfont \color{ULBblue}}

% Titre
\newcommand{\titleA}{\fontTitle{Implementation of High-Level Cryptographic}} % Titre identique au titre remis au secrétariat
\newcommand{\titleB}{\fontTitle{Protocols using a SoC Platform}} % (dans la langue de rédaction a priori)
% Sous-titre
%\newcommand{\subtitle}{\fontSubtitle{Ligne du sous-titre du mémoire}}
% Titre du diplôme
\newcommand{\diplomaA}{\fontText{Mémoire présenté en vue de l’obtention du diplôme}} % A laisser en Français
\newcommand{\diplomaB}{\fontText{d'Ingénieur Civil en informatique à finalité spécialisée}}

% Etudiant
\newcommand{\student}{\textbf{\sffamily \large Quentin Delhaye}}

% Supervision
\newcommand{\promAa}{\fontColor{Directeur}}
\newcommand{\promAb}{\fontText{Professeur Frédéric Robert}}
% \newcommand{\promBa}{\fontColor{Co-Promoteur}}
% \newcommand{\promBb}{\fontText{Professeur [Prénom Nom]}}
\newcommand{\promCa}{\fontColor{Superviseur}}
\newcommand{\promCb}{\fontText{Sébastien Rabou (Barco Silex)}}
\newcommand{\deptA}{\fontColor{Service}}
\newcommand{\deptB}{\fontText{BEAMS}}

% Année académique
\newcommand{\yearA}{\fontColor{Année académique}}
\newcommand{\yearB}{\fontText{2014 - 2015}}